\section{Literatür verisiyle yazılım uygulamaları}
\label{sec:spss-b08}

\textbf{Amaç ve veri.} \texttt{b08.xlsx} ile ortalama için $t$ güven
aralığını, \texttt{pass10} oranı için Wilson aralığını bulun
\citep{cortez2008data}.

\textbf{Ortak veri.} Üç yazılımda aynı kayıtlar ve değişkenler kullanılır. R ve Python için \texttt{companion/spss/csv/b08.csv}, Excel dosyasının aynı veri sayfasını içerir. Kodları proje kökünden çalıştırın; veri sözlüğü ve hazırlık için bkz. sayfa~\pageref{sec:spss-guide}.

\subsection{IBM SPSS uygulaması}
\textbf{Başlamadan önce.} Rehberdeki veri açma adımıyla yalnız \texttt{open-b08.sps} dosyasını çalıştırın; etiketli veri açılır. Menü yolu ve aşağıdaki tam komut dosyası alternatif uygulama yollarıdır.

\begin{spssadimlar}
\spssadim{\emph{Analyze $\to$ Descriptive Statistics $\to$ Explore} yolunda \texttt{g3} değişkenini \emph{Dependent List} alanına taşıyın. \emph{Factor List} boş kalsın.}
\spssadim{\emph{Display} için \emph{Statistics} seçin. \emph{Statistics} penceresinde \emph{Descriptives} ve yüzde 95 ortalama güven düzeyini belirleyin; \emph{Continue $\to$ OK} seçin.}
\spssadim{Çıktıdaki \emph{95\% Confidence Interval for Mean} altında \emph{Lower Bound} ve \emph{Upper Bound} satırlarını bulun. Bunlar not ortalamasının sınırlarıdır.}
\spssadim{\emph{Analyze $\to$ Descriptive Statistics $\to$ Frequencies} yolunda \texttt{pass10} tablosunu alın; 1 kategorisinin frekansını ve toplamı okuyun.}
\spssadim{Oran için Wilson aralığı, Explore içindeki ortalama aralığı değildir. Aşağıdaki tam komut dosyasını \emph{File $\to$ Open $\to$ Syntax} ile açıp \emph{Run $\to$ All} seçin; \texttt{alt} ve \texttt{ust} sınırlarını \texttt{LIST} çıktısından okuyun.}
\end{spssadimlar}

{\raggedright
\noindent\textbf{Komutların görevi.} \texttt{EXAMINE} ortalamanın $t$ aralığını verir. Sonraki \texttt{AGGREGATE} ve \texttt{COMPUTE} komutları oran için Wilson sınırlarını hesaplar; bunlar ayrı çıktılardır.\par}

\par\noindent\begin{minipage}{\linewidth}
\textbf{Uygulama kodu:} \texttt{b08}. Veri: \texttt{excel/b08.xlsx};
komut: \texttt{syntax/b08.sps} (\texttt{companion/spss} altında).

% Derleme harici .sps dosyasına bağlı değildir. Eşlikçi kopya:
% companion/spss/syntax/b08.sps
\begin{lstlisting}[style=prismSPSS]
INSERT FILE='companion/spss/syntax/open-b08.sps'
 ERROR=STOP.
EXAMINE VARIABLES=g3
 /PLOT=NONE /STATISTICS=DESCRIPTIVES
 /CINTERVAL=95.
FREQUENCIES VARIABLES=pass10.
AGGREGATE /OUTFILE=*
 /BREAK= /n=N /oran=MEAN(pass10).
COMPUTE z=1.959963984540054.
COMPUTE payda=1+z**2/n.
COMPUTE merkez=(oran+z**2/(2*n))/payda.
COMPUTE yari=z*SQRT(oran*(1-oran)/n+z**2/(4*n**2))/payda.
COMPUTE alt=merkez-yari.
COMPUTE ust=merkez+yari.
FORMATS oran alt ust (F10.4).
LIST VARIABLES=n oran alt ust.
\end{lstlisting}

\end{minipage}\par

\begin{outputguidebox}
\textbf{Beklenen kontrol değerleri (SPSS burada çalıştırılmadı).}
\emph{Descriptives} tablosunda ortalamanın yüzde 95 güven aralığı
$[\SPMatCILow;\SPMatCIHigh]$ olur. \texttt{LIST} çıktısında oran
\SPMatPHat, Wilson alt ve üst sınırları \SPWilsonLow\ ve \SPWilsonHigh'dır.
Birinci aralığın birimi not puanı, ikincisininki orandır.
Bu aralıklar öğrencilerin yüzde 95'inin notlarını kapsayan aralıklar değildir.
Model ve örnekleme varsayımları geçerliyken tekrarlanan örneklemelerden elde
edilen aralıkların kapsama özelliğiyle yorumlanırlar. İki okuldan veri
toplanmış olmasının genelleme sınırlılığı aralık hesaplayarak ortadan kalkmaz.
\end{outputguidebox}


\par\noindent\begin{minipage}{\linewidth}
\subsection{R uygulaması}
\label{sec:r-gercek-b08}

İlk satır yalnız ortalama için $t$ aralığını seçer. Sonraki hesap, SPSS ile aynı Wilson formülünü kullanır.

\begin{lstlisting}[style=prismR]
d <- read.csv("companion/spss/csv/b08.csv")
stopifnot(!anyNA(d))
print(t.test(d$g3, conf.level=.95)$conf.int)
print(table(d$pass10))
n <- nrow(d); p <- mean(d$pass10)
z <- qnorm(.975); den <- 1 + z^2/n
center <- (p + z^2/(2*n))/den
half <- z * sqrt(p*(1-p)/n + z^2/(4*n^2))/den
print(c(proportion=p, lower=center-half,
        upper=center+half))
\end{lstlisting}

\textbf{Çıktıyı okuma.} Aşağıdaki ortak karşılaştırmada verilen nicelikleri, kodun yazdırdığı etiketler ve değişken sırasıyla eşleştirin. R kodu bu ortamda çalıştırılmamıştır.
\end{minipage}\par

\par\noindent\begin{minipage}{\linewidth}
\subsection{Python uygulaması}
\label{sec:python-b08}

Ortalama aralığı ile Wilson oran aralığı ayrı hesaplanır. \texttt{wilson} alt ve üst sınırı saklar.

\begin{lstlisting}[style=prismPythonApp]
import pandas as pd
import numpy as np
from scipy import stats

d = pd.read_csv("companion/spss/csv/b08.csv")
assert not d.isna().any().any()
mean_ci = stats.ttest_1samp(d.g3, 0).confidence_interval(.95)
print(mean_ci)
print(d.pass10.value_counts().sort_index())
n, p = len(d), d.pass10.mean()
z = stats.norm.ppf(.975); den = 1 + z*z/n
center = (p + z*z/(2*n))/den
half = z*np.sqrt(p*(1-p)/n + z*z/(4*n*n))/den
wilson = (center-half, center+half)
print(p, wilson)
\end{lstlisting}

\textbf{Çıktıyı okuma.} Yazdırılan değerleri aşağıdaki ortak kontrol noktalarıyla karşılaştırın; yalnız testin $p$ değerine bakarak rapor yazmayın.
\end{minipage}\par

\subsection{Sonuçların karşılaştırılması ve ortak rapor}
\begin{outputguidebox}
Ortalama aralığı $[\SPMatCILow;\SPMatCIHigh]$, Wilson aralığı $[\SPWilsonLow;\SPWilsonHigh]$ olmalıdır. Wald veya süreklilik düzeltmeli bir aralıkla karşılaştırma yapılmamalıdır.
\end{outputguidebox}

\textbf{Raporlama örneği (kontrol değerleriyle).}
``Bağımsız gözlemler modeli altında ortalama için yüzde 95 güven aralığı $[\SPMatCILow;\SPMatCIHigh]$, eşik göstergesinin oranı için Wilson aralığı $[\SPWilsonLow;\SPWilsonHigh]$ bulunmuştur.''


\textbf{Görev.} Ortalama aralığının 10'u içerip içermediğini bulun.
Aynı veri, çift yönlü hipotez ve anlamlılık düzeyinde bunun $t$ testiyle
ilişkisini açıklayın. \texttt{.95} güven olasılığı ile \texttt{95} yüzde
gösteriminin farklı komutlarda kullanılmasına dikkat edin.